\scanpage{53}
\uline{Th.} Conditions équivalentes sur $\xi\in H^1(S,H_1)$, définissant $G=G_1^\xi$.

(i) $\xi\in\operatorname{Im}H^1(S,T_1)$ dans $H^1(S,H_1)$

(i') $\overline\xi\in\operatorname{Im}H^1(S,T'_1)$ dans $H^1(S,\Gamma_1)$ / ou $\overline\xi_0\in\operatorname{Im}H^1(S,T'_0)$ dans $H^1(S,\unsure{\Gamma'_0})$

(i'') $\overline\xi\in\operatorname{Im}H^1(S,T_1)$ dans $H^1(S,\Gamma_1)$ / ou $\overline\xi_0$

(ii) Il existe un couple de Killing dans $G$, i.e. un Borel $\supset$ Tore max.

[\deleted{(iii)} Il existe un \deleted{Tore maximal \unsure{intérieur isomorphe}} \deleted{\unsure{contenant de}} Tore maximal de $G$, \deleted{\unsure{isom. à} $T_1$} provenant \deleted{\unsure{d'un}} d'une immersion \uuline{extérieure} $T_1-G$]

(iii') Il existe une \uline{immersion intérieure} de $T_1$ dans $G$.

\margin{De plus, si $T$ est un Tore maximal, conditions équival. :
\diagram{\begin{tikzcd}[column sep=small]
\text{a) }\exists\text{ Borel}\supset T \arrow[dr,Rightarrow] & \text{b) }\exists\text{ isom }T_1\to T\ \text{extérieur}\\
&\text{c) }\exists\text{ isom }T_1\to T\ \text{intérieur}
\end{tikzcd}}}

\deleted{\uline{Dém}}
\diagram{\begin{tikzcd}[column sep=small,row sep=huge]
(iii')\arrow[r,Leftrightarrow]\arrow[d,Rightarrow,"\text{triv}"'] & (i)\arrow[r,Leftrightarrow] & (i')\arrow[r,Leftrightarrow]\arrow[d,Rightarrow,"\deleted{\unsure{triv}}"'] & \deleted{(ii)}\arrow[r,Leftrightarrow] & \deleted{\unsure{(iii')}}\\
(iii)\arrow[rr,Leftrightarrow] && (i'')\arrow[u,Rightarrow,"\text{triv}"] & &
\end{tikzcd}}
\deleted{\unsure{uniquement des lois de conjugaison}.}
\deleted{Car le centralisateur de $T_1$ en $\Gamma_1$ est $T_1$}

\uline{N.B.} Supposons que $T_1$ soit $\simeq G_m^r$ (\uline{\unsure{cas d'un anneau}} \uline{local}]. Alors \uline{$\forall$} tore $T\subset G$ qui est $\simeq T_1$, est \deleted{un tore de \unsure{Borel type}} contenu dans un Borel, \deleted{\unsure{donc}} i.e. il existe un isom \uline{intérieur} de $T_1$ sur $T$. \deleted{La \unsure{preuve}}\inserted{En effet}, \deleted{\unsure{étant local}} les \inserted{\unsure{schéma des}} isom \uline{intérieurs} de $T_1$ dans $T$ forment un sous-schéma de $\underline{\operatorname{Isom}}_S(T_1,T)\simeq(Z^{r^2})_S$, qui est fini étale sur $S$ de groupe $W$, donc \unsure{l'un est une réunion de composantes}. \uline{Or} D., \unsure{en un}

\margin{Et si $\operatorname{Pic}(S)=e$, il existe un isom de \unsure{Tch} de $G_0$ sur $G$, transformant $T_0$ en $T$.}
\margin{52}
